In this paper, we presented two ontologies: one for modelling the diversity of addresses and another for representing the temporal evolution of the geographic landmarks that compose them. Together, they support the incremental construction of a geo-historical knowledge graph from fragmentary data. However, several limitations remain and will need to be addressed in future work. First, temporal values associated with events are currently assumed to be known and precise. Since uncertainty is an inherent characteristic of historical data, future developments should explicitly account for temporal imprecision. Second, the ontology does not yet support the temporal evolution of relationships, particularly the spatial relationships linking geographic landmarks.

Following the SAMOD methodology, one remaining ontology module still needs to be developed to represent data sources and the assertions they contain. This module will enable the knowledge graph to be populated from heterogeneous information fragments while also managing conflicts between potentially contradictory sources.

Finally, to further validate the proposed ontologies, we plan to populate them with a wider variety of datasets covering Paris. These include historical cartographic sources providing \textit{snapshot}-like geometries, such as the Verniquet and Jacoubet atlases; structured datasets covering different time periods, such as the \textit{Dénominations caduques des voies de Paris}\footnote{\url{https://opendata.paris.fr/}} and the French National Address Database (BAN)\footnote{\url{https://www.data.gouv.fr/fr/datasets/base-adresse-nationale/}}; and unstructured historical sources, such as historical trade directories and textual descriptions of street histories available on Wikipedia. The resulting geo-historical knowledge graph is intended to support the development of a historical geocoding tool capable of locating historical addresses in both space and time within the city of Paris.